% Generated draft: requires mathematical and source review.
\documentclass[11pt]{article}
\usepackage[margin=25mm]{geometry}
\usepackage{fontspec}
\setmainfont{Latin Modern Roman}
\usepackage{amsmath,amssymb,mathtools,unicode-math}
\setmathfont{Latin Modern Math}
\usepackage{enumitem,needspace,xurl,hyperref}
\setlength{\parindent}{0pt}
\setlength{\parskip}{5pt}
\setlength{\emergencystretch}{3em}
\newcommand{\N}{\mathbb{N}}
\newcommand{\Z}{\mathbb{Z}}
\newcommand{\Q}{\mathbb{Q}}
\newcommand{\R}{\mathbb{R}}
\newcommand{\C}{\mathbb{C}}
\newcommand{\F}{\mathbb{F}}
\newcommand{\A}{\mathbb{A}}
\newcommand{\PP}{\mathbb{P}}
\newcommand{\E}{\mathbb{E}}
\newcommand{\eps}{\varepsilon}
\newcommand{\dd}{\,\mathrm{d}}
\begin{document}
\section*{UnsolvedMath ID: 800286}
\textbf{Source problem code:} catalogue entry
\subsection{Definitions and mathematical statement}
Let \(n\geq 1\) and let \(\Omega\subset\C^n\) be a bounded domain with \(C^\infty\) boundary. Near each \(p\in\partial\Omega\), choose a real-valued smooth defining function \(r\), so that \(\Omega=\{r<0\}\) locally and \(\dd r\neq0\) on the boundary. Pseudoconvexity means that the Levi form of \(r\) is nonnegative on the complex tangent space
\[
T^{1,0}_p\partial\Omega=\left\{v\in\C^n:\sum_{k=1}^n r_{z_k}(p)v_k=0\right\}
\]
at every boundary point.

For a nonconstant holomorphic curve germ \(\gamma:(\C,0)\to(\C^n,p)\), let \(\nu_0(\gamma-p)\) be the minimum vanishing order of its nonzero component germs. For a smooth function germ, let \(\nu_0\) denote its order of vanishing in the Taylor expansion, with value \(\infty\) for a flat germ. The D'Angelo curve type at \(p\) is
\[
\Delta_1(\partial\Omega,p)
=\sup_{\gamma}\frac{\nu_0(r\circ\gamma)}{\nu_0(\gamma-p)}.
\]
Thus finite D'Angelo type means \(\Delta_1(\partial\Omega,p)<\infty\) for every \(p\in\partial\Omega\).

To make the intended ideal-type conclusion explicit, the following is the conventional boundary \((0,1)\)-form Kohn algorithm. Let \(\mathcal E_p\) be the ring of complex-valued smooth function germs at \(p\). For an ideal \(J\subset\mathcal E_p\), write
\[
\sqrt[\mathrm R]{J}
=\left\{g\in\mathcal E_p:\text{there exist }f\in J\text{ and }m\geq1\text{ with }|g|^m\leq |f|\text{ near }p\right\}.
\]
For a differential form \(\eta\), let \(\operatorname{coeff}(\eta)\) denote its coefficient germs in local complex coordinates. Parentheses around sets of germs denote the ideal they generate. Set
\[
I_1(p)=\sqrt[\mathrm R]{\bigl(r,\operatorname{coeff}(\partial r\wedge\bar\partial r\wedge(\partial\bar\partial r)^{n-1})\bigr)}.
\]
For \(k\geq1\), let \(A_k(p)\) consist of all coefficients of
\[
\partial f_1\wedge\cdots\wedge\partial f_j\wedge\partial r\wedge\bar\partial r\wedge(\partial\bar\partial r)^{n-1-j},
\qquad 1\leq j\leq n-1,\quad f_1,\ldots,f_j\in I_k(p),
\]
and define
\[
I_{k+1}(p)=\sqrt[\mathrm R]{I_k(p)+(A_k(p))}.
\]
When \(n=1\), the displayed range for \(j\) is empty. Termination at \(p\) means that \(1\in I_N(p)\) for some finite \(N\).

\textbf{Requested statement, using this conventional interpretation of ideal type.} For every smooth bounded pseudoconvex domain \(\Omega\subset\C^n\),
\[
\left[\forall p\in\partial\Omega,\ \Delta_1(\partial\Omega,p)<\infty\right]
\quad\Longrightarrow\quad
\left[\forall p\in\partial\Omega,\ \exists N_p\geq1:\ 1\in I_{N_p}(p)\right].
\]
The source calls the conclusion ``finite type in the sense of Kohn (ideal type)'' without defining it or specifying a form degree. Its identification with the explicit algorithm above therefore requires human confirmation. No uniform bound on the termination index is asserted.
\subsection{Short English statement}
The problem asks whether every smooth bounded pseudoconvex domain of finite D'Angelo type at every boundary point also has finite Kohn ideal type there. The draft expresses the latter through termination of the conventional boundary \((0,1)\)-form Kohn algorithm; the source does not specify that convention. The registry labels the problem open, without supplying a current verification.
\subsection{Sources}
\begin{itemize}
\item Local UnsolvedMath registry, permanent ID \(800286\), exact source code \texttt{catalogue entry}, title \emph{Several Complex Variables Problem List --- Problem 86}. The supplied record attributes the problem to Buzzard, Fornaess, and Krantz, with proposed year \(1995\), and labels its status \texttt{open}.
\item Buzzard, Fornaess, and Krantz, \emph{The Several Complex Variables Problem List} (1995), Problem 86. Supplied source URL: \url{https://arxiv.org/abs/math/9509202}. The record reports access on \(2026\text{-}10\text{-}16\) and extraction by \texttt{arxiv\_latex}; these are supplied provenance claims, not independent source verification.
\item The supplied problem text includes a historical prize notice: \$50, attributed to J. McNeal, with contact \texttt{mcneal@math.princeton.edu}. No present validity of the prize or contact information is asserted.
\end{itemize}
\end{document}
